% ======================================================================
% Chapter 19: Perturbation Theory Anomalies
% (book pages 651-686; PDF pages 674-711)
% ======================================================================
\chapter{微扰论反常}

在许多例子中，我们已经看到圈修正可以对量子场论的预言产生重要影响。我们研究过这样的例子：在其中，算符的相对重要性被辐射修正所改变，以及它们所传递的相互作用的形式被改变。然而，在某些特定情形下，辐射修正可以产生更加重大的影响：它们可以破坏经典运动方程的对称性。

这类效应中最重要的涉及无质量费米子理论的手征对称性。在第~3.4 节中，我们看到无质量 Dirac 拉格朗日量具有一个增强的对称性，它与左手和右手费米子各自数目的守恒相关联。这个对称性由轴矢量流 $j^{\mu 5} = \bar{\psi}\gamma^\mu\gamma^5\psi$ 生成。经典地，
\begin{equation}
  \partial_\mu j^{\mu 5} = 0
  \label{eq:19.1}
\end{equation}
对于零质量费米子。这个运动方程不仅在自由费米子理论中成立，而且作为经典场方程，在无质量 QED 和 QCD 中也都成立。

然而，在本章中，我们将看到真实图景并非如此简单。我们将证明，在规范理论中，轴矢量流的守恒实际上与规范不变性不相容，而且规范理论中的辐射修正会提供一个非零算符，它出现在式~\eqref{eq:19.1} 的右端。轴矢流守恒的这一新方程具有若干非凡的推论，我们将在第~19.3 节和第~19.4 节中讨论它们。

\section{二维中的轴矢量流}

最终，我们将想要分析无质量 QCD 中轴矢流的流守恒方程。然而，这一讨论将涉及一些技术上的复杂性，因此我们将首先在一个计算相对简单的情形下研究破坏轴矢流守恒的物理。一个特别简单的模型问题是二维无质量 QED。

二维 QED 的拉格朗日量为
\begin{equation}
  \mathcal{L} = \bar{\psi}(i\slashed{D})\psi - \frac{1}{4}(F_{\mu\nu})^2 ,
  \label{eq:19.2}
\end{equation}
其中 $\mu, \nu = 0, 1$ 且 $D_\mu = \partial_\mu + ieA_\mu$。Dirac 矩阵必须被选取为满足 Dirac 代数
\begin{equation}
  \{\gamma^\mu, \gamma^\nu\} = 2g^{\mu\nu} .
  \label{eq:19.3}
\end{equation}
在二维中，这一组关系可以用 $2\times2$ 矩阵来表示；我们选取
\begin{equation}
  \gamma^0 = \sigma^2,\qquad \gamma^1 = i\sigma^1 .
  \label{eq:19.4}
\end{equation}
Dirac 旋量将是两分量场。

与每一个 $\gamma^\mu$ 都反对易的 Dirac 矩阵之积为
\begin{equation}
  \gamma^5 = \gamma^0\gamma^1 = \sigma^3 .
  \label{eq:19.5}
\end{equation}
于是，正如在四维中一样，存在两个可能的流，
\begin{equation}
  j^\mu = \bar{\psi}\gamma^\mu\psi;\qquad j^{\mu 5} = \bar{\psi}\gamma^\mu
    \gamma^5\psi ,
  \label{eq:19.6}
\end{equation}
如果拉格朗日量中没有质量项，则二者都守恒。

为了使守恒律非常明确，我们在这个旋量基中把费米子场 $\psi$ 的分量标记为
\begin{equation}
  \psi = \begin{pmatrix}\psi_+\\ \psi_-\end{pmatrix} .
  \label{eq:19.7}
\end{equation}
下标表示 $\gamma^5$ 本征值。然后，利用显式表示 \eqref{eq:19.4} 和 \eqref{eq:19.7}，我们可以把 \eqref{eq:19.2} 的费米子部分改写为
\begin{equation}
  \mathcal{L} = \psi^{\dagger}_+ i(D_0 + D_1)\psi_+
    + \psi^{\dagger}_- i(D_0 - D_1)\psi_- .
  \label{eq:19.8}
\end{equation}
在自由理论中，$\psi_+$ 的场方程将为
\begin{equation}
  i(\partial_0 + \partial_1)\psi_+ = 0 ;
  \label{eq:19.9}
\end{equation}
这个方程的解是以光速在一维空间中向右运动的波。因此我们将把与 $\psi_+$ 相关的粒子称为\emph{右行费米子}。类似地，与 $\psi_-$ 相关的量子为\emph{左行}的。这种区分类似于左手与右手粒子之间的区分，后者给出了 $\gamma^5$ 在四维中的物理解释。由于拉格朗日量 \eqref{eq:19.8} 不包含混合左行与右行场的项，这些场的数目流分别守恒这一点似乎是显然的。因此，
\begin{equation}
  \partial_0 N_+ = \partial_0 N_- = 0 .
  \label{eq:19.10}
\end{equation}

二维时空的一个奇特性质是，矢量与轴矢量费米子流并非彼此独立。设 $\epsilon^{\mu\nu}$ 是二维中的完全反对称符号，且 $\epsilon^{01} = +1$。则二维 Dirac 矩阵满足恒等式
\begin{equation}
  \gamma^\mu\gamma^5 = -\epsilon^{\mu\nu}\gamma_\nu .
  \label{eq:19.11}
\end{equation}
流 $j^{\mu 5}$ 与 $j^\mu$ 具有相同的关系。因此，我们可以利用已经为矢量流导出的结果来研究轴矢量流的性质。

\subsection{真空极化图}

在第~7.5 节中，我们在维数正则化下计算了 QED 的最低阶真空极化。在零质量极限下，我们在式~(7.90) 中发现，
\begin{equation}
  i\Pi^{\mu\nu}(q) = -i(q^2g^{\mu\nu} - q^\mu q^\nu)\ \frac{2e^2}{(4\pi)^{d/2}}\
    \mathrm{tr}[\mathbb{1}]\ \int_0^1 dx\ x(1-x)\ \frac{\Gamma(2-d/2)}
      {(-x(1-x)q^2)^{2-d/2}} ,
  \label{eq:19.12}
\end{equation}
其中 $\mathrm{tr}[\mathbb{1}] = 4$ 给出式~(7.88) 中给出的对 Dirac 矩阵求迹的约定。如果我们设 $\mathrm{tr}[\mathbb{1}] = 2$ 以与 \eqref{eq:19.4} 一致，然后在 \eqref{eq:19.12} 中设 $d = 2$，我们发现有限且良定义的结果
\begin{equation}
  \Pi^{\mu\nu}(q) = -\frac{e^2}{\pi}\ \frac{q^\mu q^\nu - q^2g^{\mu\nu}}
    {q^2} .
  \label{eq:19.13}
\end{equation}
注意这个表达式具有光子质量项的结构；光子获得质量
\begin{equation}
  m_\gamma^2 = \frac{e^2}{\pi} .
  \label{eq:19.14}
\end{equation}
Schwinger 证明了这个结果是精确的，并且二维 QED 的光子是自由的有质量玻色子。%
\footnote{J.~Schwinger, \emph{Phys.\ Rev.} \textbf{128}, 2425 (1962)。} 在
式~\eqref{eq:7.72} 下方的讨论中，我们指出，与 Ward 恒等式相容的真空极化振幅不可能产生光子的质量，除非它还包含 $q^2 = 0$ 处的极点。在二维中，这样的极点可以来自费米子-反费米子中间态的红外行为，而我们在 \eqref{eq:19.13} 中明确地看到这一行为。

\begin{figure}[t]
  \centering
  \includegraphics[width=0.5\textwidth]{fig19.1.png}
  \caption{在背景电磁场中 $\langle j^\mu\rangle$ 的计算。}
  \label{fig:19.1}
\end{figure}

一旦我们有了真空极化的显式表达式，我们就可以求出由背景电磁场感生的电流的期望值。这个量由图~\ref{fig:19.1} 的图产生，它给出
\begin{equation}
  \int d^2x\ e^{iq\cdot x}\ \langle j^\mu(x)\rangle
  = -\bigl( \Pi^{\mu\nu}(q) \bigr)\ A_\nu(q)
  = -\frac{e^2}{\pi}\ \frac{q^\mu q^\nu - q^2g^{\mu\nu}}{q^2}\ A_\nu(q) .
  \label{eq:19.15}
\end{equation}
其中 $A_\nu(q)$ 是背景场的 Fourier 变换。这个量显然满足流守恒关系
$q_\mu\langle j^\mu(q)\rangle = 0$。

矢量流与轴矢量流之间的恒等式 \eqref{eq:19.11} 使我们能够从 \eqref{eq:19.15} 导出 $j^{\mu 5}$ 的相应期望值。我们发现
\begin{equation}
  \langle j^{\mu 5}(q)\rangle = -\epsilon^{\mu\nu}\ \langle j_\nu(q)\rangle
  = \frac{e^2}{\pi}\ \frac{\epsilon^{\mu\nu}q_\nu q^\lambda}{q^2}\ A_\lambda(q)
  .
  \label{eq:19.16}
\end{equation}
如果轴矢量流是守恒的，这个量将满足 Ward 恒等式。但相反，我们发现
\begin{equation}
  q_\mu\langle j^{\mu 5}(q)\rangle = \frac{e^2}{2\pi}\ \epsilon^{\alpha\beta}
    F_{\alpha\beta}(q) .
  \label{eq:19.17}
\end{equation}
这是场方程的 Fourier 变换
\begin{equation}
  \partial_\mu j^{\mu 5} = \frac{e^2}{2\pi}\ \epsilon^{\alpha\beta}F_{\alpha\beta}
  .
  \label{eq:19.18}
\end{equation}
显然，作为其真空极化图反常行为的结果，轴矢量流在电磁场存在时并不守恒。

这怎么可能发生？费曼图对矢量流和轴矢量流都形式上满足 Ward 恒等式。问题必定出在真空极化图的正则化上。通过量纲分析，我们知道这个图具有形式
\begin{equation}
  \Pi^{\mu\nu}(q) = A\ q^\mu q^\nu + B\ g^{\mu\nu}q^2 .
  \label{eq:19.19}
\end{equation}
系数 $B$ 是有限积分，并且无论如何都由理论的低能结构明确地确定，因为它是 $q^2$ 处极点的留数。然而，积分 $A$ 是对数发散的，所以它的值依赖于正则化。维数正则化自动减除这个积分以令 $A = -B$；于是矢量流 Ward 恒等式得到满足。但随后我们就直接得到 \eqref{eq:19.17}。我们也可以选择把积分 $A$ 正则化使得 $A = 0$。带着这种修改重新进行上一段的步骤，我们现在发现 $q_\mu\langle j^{\mu 5}(q)\rangle = 0$，但是
\begin{equation}
  q_\mu\langle j^\mu(q)\rangle = \frac{e^2}{\pi}\ q_\mu A^\mu(q) .
  \label{eq:19.20}
\end{equation}
尽管结果 \eqref{eq:19.17} 令人不快，结果 \eqref{eq:19.20} 将是一场彻底的灾难，因为它依赖于矢势的非物理规范自由度。我们得出结论：不可能把二维 QED 正则化使得理论同时具有规范不变性并且轴矢量流守恒。要求规范不变性的代价是 \eqref{eq:19.18} 中所示的轴矢流的反常不守恒。

\subsection{轴矢量流算符方程}

为了从另一个视角理解轴矢流发生了什么，我们现在研究 $j^{\mu 5}$ 散度的算符方程。对拉格朗日量 \eqref{eq:19.2} 变分，我们得到费米子场的如下运动方程：
\begin{equation}
  i\slashed{D}\psi = 0,\qquad iD_\mu\bar{\psi}\gamma^\mu = 0 .
  \label{eq:19.21}
\end{equation}
以最直接的方式利用这些运动方程，容易得出结论 $\partial_\mu j^{\mu 5} = 0$。然而，更仔细地审视这些处理会揭示一些微妙之处，它们会改变最终结论。

轴矢量流是由费米子场构造的复合算符。在上一章中我们看到，局域算符的乘积通常是奇异的，因此我们将通过把两个费米子场放在相距 $\epsilon$ 的不同点处，然后仔细地取两个场相互靠近的极限来定义这个流。明确地说，我们定义
\begin{equation}
  j^{\mu 5}(x) = \mathrm{symm}\lim_{\epsilon\to 0}\ \Bigl[
    \bar{\psi}\Bigl( x + \frac{\epsilon}{2} \Bigr)\gamma^\mu\gamma^5\
    \exp\Bigl( -ie\int_{x-\epsilon/2}^{x+\epsilon/2} dz\cdot A(z) \Bigr)\
    \psi\Bigl( x - \frac{\epsilon}{2} \Bigr) \Bigr] .
  \label{eq:19.22}
\end{equation}
注意，因为我们把 $\bar{\psi}$ 和 $\psi$ 放在了不同点处，我们必须引入 Wilson 线 \eqref{eq:15.53} 以使该算符是局域规范不变的。为了赋予 $j^{\mu 5}$ 在 Lorentz 变换下正确的变换性质，极限 $\epsilon\to 0$ 应当对称地取，
\begin{equation}
  \mathrm{symm}\lim_{\epsilon\to 0}\Bigl\{ \frac{\epsilon^\mu}{\epsilon^2}
    \Bigr\} = 0,\qquad
  \mathrm{symm}\lim_{\epsilon\to 0}\Bigl\{ \frac{\epsilon^\mu\epsilon^\nu}
    {\epsilon^2} \Bigr\} = \frac{g^{\mu\nu}}{d} ,
  \label{eq:19.23}
\end{equation}
在现在的情形下 $d = 2$。

我们现在计算按 \eqref{eq:19.22} 定义的轴矢流的散度。利用运动方程，我们得到算符方程
\begin{equation}
  \partial_\mu j^{\mu 5} = \mathrm{symm}\lim_{\epsilon\to 0}\ \Bigl[
    \bar{\psi}\Bigl( x + \frac{\epsilon}{2} \Bigr)\gamma^\mu\gamma^5\
    \exp(\cdots)\ \psi\Bigl( x - \frac{\epsilon}{2} \Bigr) \Bigr]\
    \cdots ,
  \label{eq:19.25}
\end{equation}
其中算符乘积在短距离处的奇异项必须被仔细求值。乘积 $\bar{\psi}(x+\epsilon/2)\psi(x-\epsilon/2)$ 的短距离展开中的首项是通过用自由场传播子收缩这两个场得到的。在二维中，这给出一个奇异贡献，当与 $\gamma^\mu\gamma^5$ 求迹时，它在 $\epsilon\to 0$ 极限下产生非零结果。这就是反常的算符形式：
\begin{equation}
  \partial_\mu j^{\mu 5} = \frac{e^2}{2\pi}\ \epsilon^{\alpha\beta}F_{\alpha\beta}
  .
  \label{eq:19.29}
\end{equation}

\subsection{一个费米子数不守恒的例子}

轴矢流的反常不守恒有一个戏剧性的推论：它意味着左手与右手费米子的数目可以被背景规范场改变。为了在一个简单例子中看到这一点，我们考虑长度为 $L$ 的圆上的二维 QED，费米子满足周期性边界条件。我们把这个理论耦合到一个恒定的背景电场，它对应于与时间相关的 $A^1$ 分量。

系统的哈密顿量为
\begin{equation}
  H = \int dx\ \Bigl[ \psi^{\dagger}_+(i\partial_1 - eA^1)\psi_+
    + \psi^{\dagger}_-(-i\partial_1 + eA^1)\psi_- \Bigr] ,
  \label{eq:19.34}
\end{equation}
其中 $\psi_{\pm}$ 是右行和左行分量。对于恒定的 $A^1$ 场，容易对角化这个哈密顿量。协变导数的本征态是波函数
\begin{equation}
  e^{ik_n x},\qquad \text{with } k_n = \frac{2\pi}{L}n,\qquad n = -\infty,
    \ldots, \infty ,
  \label{eq:19.35}
\end{equation}
以满足周期性边界条件。则 $H$ 的单粒子本征态具有能量
\begin{equation}
  E^+_n = +(k_n - eA^1),\qquad E^-_n = -(k_n - eA^1) .
  \label{eq:19.36}
\end{equation}
每种类型的费米子都有一列无限延伸的等间距能级。为了求 $H$ 的基态，我们填满负能级，并把这些被填充态中产生的空穴解释为反粒子。

\begin{figure}[t]
  \centering
  \includegraphics[width=0.55\textwidth]{fig19.2.png}
  \caption{一维 QED 的哈密顿量 $H$ 的真空态因背景 $A^1$ 场的绝热改变而受到的影响。}
  \label{fig:19.2}
\end{figure}

现在绝热地改变 $A^1$ 的值。费米子能级按照关系 \eqref{eq:19.36} 缓慢移动。如果 $A^1$ 改变有限量
\begin{equation}
  \Delta A^1 = \frac{2\pi}{eL} ,
  \label{eq:19.37}
\end{equation}
它使 Wilson 圈式~(19.32) 回到其原值，$H$ 的谱回到其原始形式。在这个过程中，$\psi_+$ 的每一个能级下移到下一个位置，$\psi_-$ 的每一个能级上移到下一个位置，如图~\ref{fig:19.2} 所示。能级的占据数在这个绝热过程中应当保持不变。因此，令人惊奇的是，一个右行费米子从真空中消失，同时出现一个额外的左行费米子。与此同时，
\begin{equation}
  \int d^2x\ \frac{e}{2\pi}\ \epsilon^{\mu\nu}F_{\mu\nu}
  = \int dt\,dx\ \frac{e}{\pi}\ \partial_0 A^1
  = \frac{e}{\pi}\ \Delta A^1\,L
  = \frac{e}{\pi}\cdot\frac{2\pi}{eL}\cdot L = 2 ,
  \label{eq:19.38}
\end{equation}
其中我们在最后一行代入了 \eqref{eq:19.37}。因此反常不守恒方程 \eqref{eq:19.18} 的积分形式确实得到满足：
\begin{equation}
  N_R - N_L = \int d^2x\ \frac{e}{2\pi}\ \epsilon^{\mu\nu}F_{\mu\nu} .
  \label{eq:19.39}
\end{equation}
即使在这个简单例子中，我们也看到在分析手征守恒律时不可能回避紫外正则化问题。右行费米子消失，左行费米子从费米子谱的深处 $E\to -\infty$ 出现。在计算各自费米子数的变化时，我们已经假设真空不能改变它在很大负能量处所含的电荷。这个方案是规范不变的，但它导致轴矢量流的不守恒。

\section{四维中的轴矢量流}

我们刚才为二维轴矢流给出的所有推导在四维中都有对应物。在式~\eqref{eq:3.40} 中，我们证明，在无质量费米子的情形下，四维 Dirac 方程干净地分裂为左手与右手费米子的独立方程。如果我们把 Dirac 方程耦合到规范场，我们就把导数替换为协变导数。这似乎不影响两个螺旋度分量之间明显的分离。因此似乎清楚的是，矢量流和轴矢量流都应保持守恒。然而，在我们刚刚完成对二维情形的分析之后，我们知道不应把这些守恒律视为理所当然。我们现在将对四维中的轴矢量守恒律做更仔细的分析。

\subsection{轴矢量流算符方程}

我们从无质量四维 QED 的情形开始。我们在上一节为二维轴矢流守恒律给出的三个论证中，算符推导最容易推广。费米子场方程 \eqref{eq:19.21} 在四维情形下是相同的。我们可以再次采用轴矢量流的规范不变定义 \eqref{eq:19.22}。当我们取这个流的散度时，所有通向式~\eqref{eq:19.25} 的处理仍然正确。

\begin{figure}[t]
  \centering
  \includegraphics[width=0.5\textwidth]{fig19.3.png}
  \caption{在背景规范场存在下 $\bar{\psi}(y)\psi(z)$ 的展开。}
  \label{fig:19.3}
\end{figure}

从这一点开始，我们必须计算两个费米子场的算符乘积在 $\epsilon\to 0$ 极限下的奇异项。正如在二维中一样，首项由使用自由场传播子收缩这两个算符给出。这个贡献给出
\begin{equation}
  \bigl\langle \bar{\psi}(y)\psi(z) \bigr\rangle \to -i\slashed{\partial}\
    \frac{1}{2\pi^2(y-z)^2} = -i\gamma^\alpha\ \frac{(y-z)_\alpha}
      {2\pi^2(y-z)^4} .
  \label{eq:19.40}
\end{equation}
这在 $(y-z)\to 0$ 时高度奇异，但当与 $\gamma^\mu\gamma^5$ 求迹时给出零。为了找到非零结果，我们必须考虑算符乘积展开中的更高阶项。

在非零背景规范场中，费米子场的收缩由图~\ref{fig:19.3} 所示的一系列图给出。我们已经在 \eqref{eq:19.40} 中计算了这个级数的首项。高阶项在 $(y-z)\to 0$ 时给出不太奇异的项。级数中的第二项由下式给出
\begin{equation}
  \text{(second term)} = \int \frac{d^4k}{(2\pi)^4}\ \frac{d^4p}{(2\pi)^4}\
    e^{ik\cdot y}\ e^{-ip\cdot z}\ \frac{(-ie\slashed{A}(p))}{k^2(k+p)^2}
    \ \cdots ,
  \label{eq:19.41}
\end{equation}
这个贡献导致
\begin{equation}
  \Bigl\langle \bar{\psi}\Bigl( x+\frac{\epsilon}{2} \Bigr)\gamma^\mu\gamma^5
    \psi\Bigl( x-\frac{\epsilon}{2} \Bigr) \Bigr\rangle
  \sim -4e\ \epsilon^{\mu\alpha\beta\gamma}\ \partial_\beta
    \partial_\gamma A_\alpha(x)\ \Bigl( \frac{1}{\epsilon^2} \Bigr)\ \cdots .
  \label{eq:19.42}
\end{equation}
为了求值 $\epsilon\to 0$ 极限，我们可以对大 $k$ 展开被积函数。于是
\begin{equation}
  \Bigl\langle \bar{\psi}\Bigl( x+\frac{\epsilon}{2} \Bigr)\gamma^\mu\gamma^5
    \psi\Bigl( x-\frac{\epsilon}{2} \Bigr) \Bigr\rangle
  \to 4e\ \epsilon^{\mu\nu\rho\sigma}\ \partial_\rho A_\sigma(x)\ \Bigl(
    \frac{1}{4\pi^2}\log\frac{1}{\epsilon^2} \Bigr) .
  \label{eq:19.43}
\end{equation}
把这个表达式代入 \eqref{eq:19.25}，我们发现
\begin{equation}
  \partial_\mu j^{\mu 5} = \mathrm{symm}\lim_{\epsilon\to 0}\ \Bigl[
    \frac{e^2}{8\pi^2}\ \epsilon^{\mu\nu\rho\sigma}F_{\mu\nu}F_{\rho\sigma}
    \Bigr] .
  \label{eq:19.44}
\end{equation}
现在在四维中取对称极限 $\epsilon\to 0$。我们发现
\begin{equation}
  \partial_\mu j^{\mu 5} = \frac{e^2}{16\pi^2}\ \epsilon^{\mu\nu\rho\sigma}
    F_{\mu\nu}F_{\rho\sigma} .
  \label{eq:19.45}
\end{equation}
这就是\emph{Adler-Bell-Jackiw 反常}方程。右端是一个新算符，它作为辐射修正的结果出现在轴矢流的散度中。

\subsection{泛函积分的手征变换}

用另一种方式、利用泛函积分来推导轴反常是很有意思的。这个推导特别清晰，并且容易推广到其他理论。考虑无质量 QED 的泛函积分，
\begin{equation}
  Z = \int \mathcal{D}\bar{\psi}\,\mathcal{D}\psi\,\mathcal{D}A\ \exp\Bigl[
    i\int d^4x\ \bigl( \bar{\psi}i\slashed{D}\psi
      - \frac{1}{4}(F_{\mu\nu})^2 \bigr) \Bigr] .
  \label{eq:19.61}
\end{equation}
我们现在对费米子场做手征变换，
\begin{equation}
  \psi(x) \to e^{i\alpha(x)\gamma^5}\psi(x),\qquad
  \bar{\psi}(x) \to \bar{\psi}(x)e^{i\alpha(x)\gamma^5} .
  \label{eq:19.62}
\end{equation}
拉格朗日量的变换为
\begin{equation}
  \bar{\psi}i\slashed{D}\psi \to \bar{\psi}i\slashed{D}\psi
    - \bar{\psi}\slashed{D}\alpha\,\gamma^5\psi
    = \bar{\psi}i\slashed{D}\psi + \alpha\,\partial_\mu j^{\mu 5} .
  \label{eq:19.63}
\end{equation}
如果费米子测度 $\mathcal{D}\bar{\psi}\,\mathcal{D}\psi$ 在手征变换下是不变的，那么对变换后的积分关于 $\alpha(x)$ 变分将给出 $\partial_\mu j^{\mu 5} = 0$，即朴素守恒律。然而，我们现在将证明这个测度并不是不变的，而且变换的 Jacobi 行列式恰好产生反常项。

为了计算 Jacobi 行列式，我们在 $i\slashed{D}$ 的本征态基中展开费米子场。定义 $i\slashed{D}$ 的右、左本征矢为
\begin{equation}
  i\slashed{D}\phi_m = \lambda_m\phi_m,\qquad
  \phi^{\dagger}_m\,i\slashed{D} = \lambda^{*}_m\phi^{\dagger}_m .
  \label{eq:19.64}
\end{equation}
对于零背景 $A^\mu$ 场，这些本征态是具有确定动量的 Dirac 波函数；本征值满足
\begin{equation}
  \lambda^2_m = k^2 = (k^0)^2 - (\mathbf{k})^2 .
  \label{eq:19.65}
\end{equation}
对于固定的 $A^\mu$ 场，这也是大 $k$ 时本征值的渐近形式。这些本征函数给我们一个基，我们可以用它来展开 $\psi$ 和 $\bar{\psi}$：
\begin{equation}
  \psi(x) = \sum_m a_m\phi_m(x),\qquad
  \bar{\psi}(x) = \sum_m \bar{a}_m\phi^{\dagger}_m(x) ,
  \label{eq:19.66}
\end{equation}
其中 $a_m$、$\bar{a}_m$ 是乘以 c 数本征函数 \eqref{eq:19.64} 的反对易系数。关于 $\psi$、$\bar{\psi}$ 的泛函测度于是可以定义为
\begin{equation}
  \mathcal{D}\bar{\psi}\,\mathcal{D}\psi = \prod_m d\bar{a}_m\,da_m ,
  \label{eq:19.67}
\end{equation}
而关于 $\psi'$、$\bar{\psi}'$ 的泛函测度可以以相同方式定义。

如果 $\psi'(x) = (1+i\alpha(x)\gamma^5)\psi(x)$，$\psi$ 和 $\psi'$ 的展开系数由无穷小线性变换 $(1+C)$ 相关联，计算如下：
\begin{equation}
  a'_m = \sum_n \int d^4x\ \phi^{\dagger}_m(x)\ (1+i\alpha(x)\gamma^5)\
    \phi_n(x)\ a_n = \sum_n \bigl( \delta_{mn} + C_{mn} \bigr)a_n .
  \label{eq:19.68}
\end{equation}
则
\begin{equation}
  \prod_m da'_m = \prod_m da_m\ \cdot\ J ,
  \label{eq:19.69}
\end{equation}
其中 $J$ 是变换 $(1+C)$ 的 Jacobi 行列式。$J$ 的逆出现于 \eqref{eq:19.69} 中，这是费米子积分规则 \eqref{eq:9.63} 或式~(9.69) 的结果。为了求值 $J$，我们写出
\begin{equation}
  J = \det(1+C) = \exp\bigl[ \mathrm{tr}\log(1+C) \bigr]
  = \exp\Bigl\{ \sum_n C_{nn} + \cdots \Bigr\} ,
  \label{eq:19.70}
\end{equation}
并且我们可以忽略最后一行中的高阶项，因为 $C$ 是无穷小的。因此，
\begin{equation}
  \log J = i\int d^4x\ \alpha(x)\ \sum_n \phi^{\dagger}_n(x)\gamma^5\phi_n(x)
  .
  \label{eq:19.71}
\end{equation}
$\alpha(x)$ 的系数看起来像 $\mathrm{tr}[\gamma^5] = 0$。然而，我们必须以规范不变的方式对 $n$ 的本征态求和做正则化。自然的选择是
\begin{equation}
  \sum_n \phi^{\dagger}_n(x)\gamma^5\phi_n(x)
  = \lim_{M\to\infty}\ \sum_n \phi^{\dagger}_n(x)\gamma^5\phi_n(x)\,
    e^{\lambda_n^2/M^2} .
  \label{eq:19.72}
\end{equation}
如 \eqref{eq:19.65} 所示，经过 Wick 转动后，$\lambda^2$ 的符号在大动量处将为负；因此，收敛因子指数中的符号给定是正确的。我们可以把 \eqref{eq:19.72} 写成算符形式
\begin{equation}
  = \lim_{M\to\infty}\ \sum_n \phi^{\dagger}_n(x)\gamma^5\
    e^{(i\slashed{D})^2/M^2}\phi_n(x)
  = \lim_{M\to\infty}\ \mathrm{tr}\Bigl[ \gamma^5\ \langle x|\,
    e^{(i\slashed{D})^2/M^2}\,|x\rangle \Bigr] ,
  \label{eq:19.73}
\end{equation}
其中在第二行中，我们对 Dirac 指标求迹。

为了求值 \eqref{eq:19.73}，我们按照 \eqref{eq:16.107} 改写 $(i\slashed{D})^2$。在我们当前的约定下，这个方程读作
\begin{equation}
  (i\slashed{D})^2 = -D^2 + \frac{e}{2}\sigma^{\mu\nu}F_{\mu\nu} .
  \label{eq:19.74}
\end{equation}
由于我们正在取极限 $M\to\infty$，我们可以把注意力集中在谱的渐近部分，那里动量 $k$ 很大，我们可以按规范场的幂次展开。为了得到与 $\gamma^5$ 的非零迹，我们必须从指数中取出四个 Dirac 矩阵。首项由把指数展开到 $(\sigma\cdot F)^2$ 阶、然后在所有其他项中忽略背景 $A^\mu$ 场给出。这给出
\begin{equation}
  \lim_{M\to\infty}\ \mathrm{tr}\Bigl[ \gamma^5\ \langle x|\,
    e^{(-D^2+(e/2)\sigma\cdot F)/M^2}\,|x\rangle \Bigr]
  = \lim_{M\to\infty}\ \mathrm{tr}\Bigl[ \gamma^5\ \frac{e^2}{8M^4}\,
    (\sigma\cdot F)^2\ \langle x|\,e^{-\partial^2/M^2}\,|x\rangle \Bigr] .
  \label{eq:19.75}
\end{equation}
\eqref{eq:19.75} 中的矩阵元可以通过 Wick 转动求值：
\begin{equation}
  \langle x|\,e^{-\partial^2/M^2}\,|x\rangle
  = \int \frac{d^4k_E}{(2\pi)^4}\ e^{-k_E^2/M^2}
  = \frac{M^4}{16\pi^2} .
  \label{eq:19.76}
\end{equation}
则 \eqref{eq:19.75} 化简为
\begin{equation}
  \lim_{M\to\infty}\ \frac{-ie^2}{8\cdot16\pi^2}\ \mathrm{tr}\bigl[
    \gamma^5\sigma^{\mu\nu}\sigma^{\rho\sigma} \bigr]\ F_{\mu\nu}F_{\rho\sigma}
  = \frac{e^2}{32\pi^2}\ \epsilon^{\mu\nu\rho\sigma}F_{\mu\nu}F_{\rho\sigma}
  .
  \label{eq:19.77}
\end{equation}
因此，
\begin{equation}
  J = \exp\Bigl[ i\int d^4x\ \alpha(x)\ \frac{e^2}{32\pi^2}\,
    \epsilon^{\mu\nu\rho\sigma}F_{\mu\nu}F_{\rho\sigma}(x) \Bigr] .
  \label{eq:19.78}
\end{equation}

总之，我们发现，在变量替换 \eqref{eq:19.62} 之后，泛函积分 \eqref{eq:19.61} 采取形式
\begin{equation}
  Z = \int \mathcal{D}\bar{\psi}\,\mathcal{D}\psi\,\mathcal{D}A\ \exp\Bigl[
    i\int d^4x\ \Bigl( \mathcal{L} + \alpha(x)\Bigl( \partial_\mu j^{\mu 5}
      + \frac{e^2}{32\pi^2}\epsilon^{\mu\nu\rho\sigma}F_{\mu\nu}F_{\rho\sigma}
      \Bigr) \Bigr) \Bigr] .
  \label{eq:19.79}
\end{equation}
对指数关于 $\alpha(x)$ 变分，我们精确地得到 Adler-Bell-Jackiw 反常方程。

轴矢量反常的这个推导特别有意思，因为它可以很容易地推广到任何偶数维。泛函推导总是为反常方程的右端挑出由规范场构造的赝标量算符，它具有与流的散度相同的量纲 $d$。在二维中，这个推导立即导致 \eqref{eq:19.18}。只要 $d$ 是偶数，我们总可以通过取所有 Dirac 矩阵之积来构造一个与它们全部反对易的矩阵 $\gamma^5$。于是，泛函推导直接导致结果
\begin{equation}
  \partial_\mu j^{\mu 5} = \frac{1}{(2\pi)^n}\ \epsilon^{\mu_1\cdots\mu_{2n}}\
    F_{\mu_1\mu_2}\cdots F_{\mu_{2n-1}\mu_{2n}} ,
  \label{eq:19.80}
\end{equation}
其中 $n = d/2$。

在上一节的末尾，我们论证了轴矢量反常导致二维系统中具有宏观电场的费米子电荷的全局不守恒。以同样的方式，四维反常方程导致在 \eqref{eq:19.45} 右端非零的背景场中左手与右手费米子数目的全局不守恒。这些是电场与磁场平行的场位形。在习题~19.1 中，我们研究这种类型的一个简单情形下的四维无质量费米子例子，并证明费米子数确实被违反，其方式类似于我们在第~19.1 节末尾看到的，且与 Adler-Bell-Jackiw 反常一致。

\section{QCD 中的 Goldstone 玻色子与手征对称性}

Adler-Bell-Jackiw 反常对 QCD 有若干重要的影响。为了描述这些影响，我们必须首先系统地讨论 QCD 的手征对称性。在这个讨论中，我们将忽略除最轻的 $u$ 和 $d$ 夸克之外的一切夸克。在许多对强相互作用低能结构的分析中，人们也把 $s$ 夸克当作轻夸克处理；这给出的结果自然推广我们下面将找到的结果。

QCD 拉格朗日量的费米子部分为
\begin{equation}
  \mathcal{L} = \bar{u}i\slashed{D}u + \bar{d}i\slashed{D}d
    - m_u\bar{u}u - m_d\bar{d}d .
  \label{eq:19.81}
\end{equation}
如果 $u$ 和 $d$ 夸克非常轻，最后两项很小，可以忽略。让我们研究作这个近似的推论。如果我们忽略 $u$ 和 $d$ 质量，拉格朗日量 \eqref{eq:19.81} 当然具有同位旋对称性，即混合 $u$ 和 $d$ 场的 $SU(2)$ 幺正变换的对称性。然而，由于无质量费米子的经典拉格朗日量不包含左手与右手夸克之间的耦合，这个拉格朗日量实际上在独立的幺正变换下是对称的
\begin{equation}
  q_L \to U_L q_L,\qquad q_R \to U_R q_R .
  \label{eq:19.82}
\end{equation}
把这些变换的 $U(1)$ 和 $SU(2)$ 部分分离开是有用的；于是经典的、无质量 QCD 拉格朗日量的对称群为 $SU(2)\times SU(2)\times U(1)\times U(1)$。设 $q$ 表示夸克双重态，其手征分量为
\begin{equation}
  q = \begin{pmatrix}u\\ d\end{pmatrix},\qquad
  q_L = \frac{1-\gamma^5}{2}q,\qquad q_R = \frac{1+\gamma^5}{2}q .
  \label{eq:19.83}
\end{equation}
则我们可以把这些对称性相关联的流写为
\begin{equation}
  j^\mu_L = \bar{q}_L\gamma^\mu q_L,\qquad j^\mu_R = \bar{q}_R\gamma^\mu q_R,
  \qquad j^{\mu a}_L = \bar{q}_L\gamma^\mu\tau^a q_L,\qquad
  j^{\mu a}_R = \bar{q}_R\gamma^\mu\tau^a q_R ,
  \label{eq:19.84}
\end{equation}
其中 $\tau^a = \sigma^a/2$ 表示 $SU(2)$ 的生成元。左手与右手流之和给出重子数与同位旋流
\begin{equation}
  j^\mu = j^\mu_L + j^\mu_R,\qquad j^{\mu a} = j^{\mu a}_L + j^{\mu a}_R .
  \label{eq:19.85}
\end{equation}
相应的对称性是 $U_L = U_R$ 的变换 \eqref{eq:19.82}。流 \eqref{eq:19.84} 之差给出相应的轴矢量流 $j^{\mu a}_5$：
\begin{equation}
  j^{\mu a}_5 = j^{\mu a}_R - j^{\mu a}_L .
  \label{eq:19.86}
\end{equation}
在接下来的讨论中，我们将通过假设这些流的经典守恒律不被反常破坏来推导关于强相互作用的结论。我们下面将证明，这个假设对同位旋三重态流 $j^{\mu a}_5$ 是正确的，但对 $j^{\mu 5}$ 则不然。

矢量 $SU(2)\times U(1)$ 变换是强相互作用的明显对称性。然而，轴变换不是基态的对称性：它们被自发破缺。算符 $\bar{q}q$ 的真空期望值非零，
\begin{equation}
  \langle 0|\bar{u}u|0\rangle = \langle 0|\bar{d}d|0\rangle \neq 0 ,
  \label{eq:19.87}
\end{equation}
它破坏轴对称性。与 $SU(2)$ 轴对称性的自发破缺相关联的 Goldstone 玻色子是 $\pi$ 介子。流 $j^{\mu a}_5$ 从真空中产生一个 $\pi$ 介子：
\begin{equation}
  \langle 0|j^{\mu a}_5(x)|\pi^b(p)\rangle = -ip^\mu f_\pi\,\delta^{ab}\,
    e^{-ip\cdot x} ,
  \label{eq:19.88}
\end{equation}
其中 $f_\pi$ 是 $\pi$ 介子衰变常数。

手征对称性的自发破缺有许多推论。特别是，它蕴含 $\pi$ 介子质量、夸克质量与轴矢流矩阵元之间的关系。$\pi$ 介子是破缺对称性的 Goldstone 玻色子，它的相互作用由从这个事实得出的低能定理决定。为了说明这些思想，让我们考虑核子态之间的轴矢流顶角。我们把轴同位旋流的核子矩阵元写为
\begin{equation}
  \langle N(p')|j^{\mu a}_5|N(p)\rangle = \bar{u}(p')\ \Bigl[
    \gamma^\mu\gamma^5 g_A(q^2) + \frac{q^\mu}{2m_N}\gamma^5\ g_P(q^2)
    \Bigr]\ \tau^a\,u(p) ,
  \label{eq:19.97}
\end{equation}
其中 $q = p'-p$。轴矢流的守恒蕴含
\begin{equation}
  q_\mu\langle N(p')|j^{\mu a}_5|N(p)\rangle = 0 .
  \label{eq:19.98}
\end{equation}
这个方程蕴含，除非流矩阵元包含 $q^2$ 中的极点，否则 $g_A$ 必为零。这样的极点将意味着存在一个物理上的无质量粒子，但幸运的是，恰好有一个可用的——无质量 $\pi$ 介子。流产生一个随后被核子吸收的 $\pi$ 介子的过程确实导致形状因子中的极点，如图~\ref{fig:19.6}(b) 所示。

\begin{figure}[t]
  \centering
  \includegraphics[width=0.55\textwidth]{fig19.6.png}
  \caption{对轴矢流核子矩阵元有贡献的图：(a) 直接耦合；(b) $\pi$ 介子极点贡献。}
  \label{fig:19.6}
\end{figure}

现在让我们计算这个极点项并用它来确定 $g_A$。低能 $\pi$ 介子-核子相互作用按惯例由拉格朗日量参数化
\begin{equation}
  \Delta\mathcal{L} = ig_{\pi NN}\,\bar{N}\gamma^5\tau^a N\,\pi^a .
  \label{eq:19.99}
\end{equation}
流 $j^{\mu a}_5$ 产生 $\pi$ 介子的振幅由 \eqref{eq:19.88} 给出。则图~\ref{fig:19.6}(b) 对电流顶角的贡献为
\begin{equation}
  \frac{1}{q^2}\ (if_\pi q^\mu)\ g_{\pi NN}\ \bar{N}\gamma^5\tau^a N .
  \label{eq:19.100}
\end{equation}
因此，
\begin{equation}
  g_P(q^2) = \frac{2m_N f_\pi g_{\pi NN}}{q^2} .
  \label{eq:19.101}
\end{equation}
利用守恒条件 \eqref{eq:19.98}，我们发现 $g_A$ 由 $\pi$ 介子衰变常数、核子质量和 $\pi$ 介子-核子耦合常数的组合给出：
\begin{equation}
  g_A = \frac{f_\pi g_{\pi NN}}{m_N} .
  \label{eq:19.102}
\end{equation}
这个奇怪的恒等式被称为\emph{Goldberger-Treiman 关系}，实验上以 5\% 的精度得到满足。

把 $\pi$ 介子认定为自发破缺手征对称性的 Goldstone 玻色子，导致对流矩阵元和 $\pi$ 介子散射振幅的许多其他预言。特别是，低能下 $\pi$-$\pi$ 和 $\pi$ 介子-核子散射振幅的首项可以通过与刚才给出的论证类似的论证直接用 $f_\pi$ 计算。%
\footnote{自发破缺手征对称性的详细推论在 Georgi (1984) 中以非常清晰的方式被导出。}

\subsection{手征流的反常}

到目前为止，我们根据经典流守恒方程讨论了 QCD 的手征对称性。我们现在必须问，这些方程是否受 Adler-Bell-Jackiw 反常的影响，以及这种改变的推论是什么。

\begin{figure}[t]
  \centering
  \includegraphics[width=0.5\textwidth]{fig19.7.png}
  \caption{在 QCD 中导致手征流的轴矢量反常的图。}
  \label{fig:19.7}
\end{figure}

首先，我们研究由于夸克流耦合到 QCD 胶子场而导致的手征守恒律的改变。上一节给出的论证在无质量费米子耦合到非阿贝尔规范场的情形下同样成立，因此我们预期轴矢量流将从图~\ref{fig:19.7} 所示的图接收反常贡献。反常方程应当是阿贝尔结果，并补充适当的群论因子。此外，由于轴矢流是规范不变的，反常也必须是规范不变的。也就是说，它必须包含完整的非阿贝尔场强，包括其非线性项。这些项实际上已经包含在第~19.2 节末尾给出的反常的泛函推导中。

对于 QCD 的轴矢流（写于 \eqref{eq:19.86} 中），我们可以从图~\ref{fig:19.7} 的图中读出 Adler-Bell-Jackiw 反常的群论因子。对于轴同位旋流，
\begin{equation}
  \partial_\mu j^{\mu a}_5 = \frac{g^2}{32\pi^2}\ \epsilon^{\mu\nu\rho\sigma}\
    F^c_{\mu\nu}F^d_{\rho\sigma}\ \mathrm{tr}\bigl[ \tau^a\{t^c, t^d\} \bigr],
  \label{eq:19.103}
\end{equation}
其中 $F^c_{\mu\nu}$ 是胶子场强，$\tau^a$ 是同位旋矩阵，$t^c$ 是色矩阵，并且迹在颜色和味上求。在这种情形下，我们发现
\begin{equation}
  \mathrm{tr}\bigl[ \tau^a t^c t^d \bigr] = \mathrm{tr}[\tau^a]\,
    \mathrm{tr}[t^c t^d] = 0 ,
  \label{eq:19.104}
\end{equation}
因为单个 $\tau^a$ 的迹为零。因此轴同位旋流的守恒不受 QCD 的 Adler-Bell-Jackiw 反常的影响。然而，在同位旋单态轴矢流的情形下，矩阵 $\tau^a$ 被味上的矩阵 1 替换，我们发现
\begin{equation}
  \partial_\mu j^{\mu 5} = n_f\ \frac{g^2}{32\pi^2}\ \epsilon^{\mu\nu\rho\sigma}
    F^c_{\mu\nu}F^c_{\rho\sigma} ,
  \label{eq:19.105}
\end{equation}
其中 $n_f$ 是味的数目；在我们当前的模型中 $n_f = 2$。

因此，同位旋单态轴矢流在 QCD 中实际上并不守恒。这个流的散度等于一个在强子态之间具有非平庸矩阵元的胶子算符。关于这个算符的效应仍存在一些微妙的问题。特别是，可以证明（正如我们在式~(19.31) 中对二维轴反常看到的），\eqref{eq:19.105} 的右端是一个全散度。然而，再次与我们在二维中的经验一致，存在物理上合理的场位形，其中这一项的四维积分取非零值。这个主题在第~22.3 节末尾被进一步讨论。无论如何，式~\eqref{eq:19.105} 确实蕴含 QCD 没有同位旋单态轴对称性，也没有相关联的 Goldstone 玻色子。这个方程解释了为什么强相互作用不包含质量与 $\pi$ 介子相当的轻同位旋单态赝标量介子。

尽管轴同位旋流没有来自 QCD 相互作用的轴反常，但它们确实具有与夸克耦合到电磁学相关联的反常。再次参考图~\ref{fig:19.7} 的图，我们看到轴同位旋流的电磁反常由下式给出
\begin{equation}
  \partial_\mu j^{\mu a}_5 = \frac{e^2}{32\pi^2}\ \epsilon^{\mu\nu\rho\sigma}\
    F_{\mu\nu}F_{\rho\sigma}\ \mathrm{tr}\bigl[ \tau^a Q^2 \bigr] ,
  \label{eq:19.106}
\end{equation}
其中 $F_{\mu\nu}$ 是电磁场强，$Q$ 是夸克电荷矩阵，
\begin{equation}
  Q = \begin{pmatrix}\frac{2}{3}&0\\ 0&-\frac{1}{3}\end{pmatrix},
  \label{eq:19.107}
\end{equation}
并且迹再次在味和颜色上求。由于迹中的矩阵不依赖于颜色，色求和只给出因子 3。味迹只在 $a = 3$ 时非零；在那种情形下，电磁反常为
\begin{equation}
  \partial_\mu j^{\mu 3}_5 = \frac{e^2}{16\pi^2}\ \epsilon^{\mu\nu\rho\sigma}\
    F_{\mu\nu}F_{\rho\sigma} .
  \label{eq:19.108}
\end{equation}
因为流 $j^{\mu 3}_5$ 湮灭一个 $\pi^0$ 介子，式~\eqref{eq:19.108} 表明轴矢量反常对衰变 $\pi^0\to 2\gamma$ 的矩阵元有贡献。我们现在将证明，事实上，它给出这个振幅的首项贡献。

再次，我们在无质量 $u$ 和 $d$ 夸克的极限下工作，从而手征对称性除反常效应外是精确的。考虑轴矢流在真空与双光子态之间的矩阵元：
\begin{equation}
  \langle \gamma(p)\gamma(k)|j^{\mu 5}(q)|0\rangle = M^{\mu\nu\lambda}
    \epsilon_\nu(p)\epsilon_\lambda(k) .
  \label{eq:19.109}
\end{equation}
这就是我们在第~19.2 节的 QED 微扰论中研究过的同一个矩阵元式~(19.46)。然而现在，我们将通过把它按形状因子展开来研究这个矩阵元的一般性质。一般地，振幅可以通过写出所有可能的张量结构并应用由 $(\mu,\nu)$ 与 $(k,\lambda)$ 互换下的对称性以及 QED Ward 恒等式式~(19.51) 推出的限制来分解。这留下三个可能的结构：
\begin{equation}
  \begin{aligned}
    M^{\mu\nu\lambda} &= \epsilon^{\mu\nu\lambda\alpha} p_\alpha k^\nu\,
      M_1 + \bigl( \epsilon^{\mu\nu\alpha\beta}k^\lambda - \epsilon^{\mu\lambda
        \alpha\beta}p^\nu \bigr)\,k_\alpha p_\beta\,M_2\\
    &\quad + \Bigl[ \bigl( \epsilon^{\mu\nu\alpha\beta}p^\lambda
      - \epsilon^{\mu\lambda\alpha\beta}k^\nu \bigr)\,k_\alpha p_\beta
      \cdot k \Bigr]\,M_3 .
  \end{aligned}
  \label{eq:19.110}
\end{equation}
第二项由于在壳条件 $p^2 = k^2 = 0$ 而满足式~(19.51)。

现在把 \eqref{eq:19.110} 与 $(iq_\mu)$ 收缩以取轴矢量流的散度。我们发现
\begin{equation}
  iq_\mu M^{\mu\nu\lambda} = iq^2\ \epsilon^{\nu\lambda\alpha\beta}\ p_\alpha
    k_\beta\ M_1 - i\ \epsilon^{\nu\lambda\alpha\beta}\ q_\alpha(p-k)_\beta\,
    k\cdot p\ M_3 ;
  \label{eq:19.111}
\end{equation}
其他项自动给出零。利用 $q = p + k$ 和 $q^2 = 2p\cdot k$，我们可以把它化简为
\begin{equation}
  iq_\mu M^{\mu\nu\lambda} = iq^2\ \epsilon^{\nu\lambda\alpha\beta}\ p_\alpha
    k_\beta\ \bigl( M_1 + M_3 \bigr) .
  \label{eq:19.112}
\end{equation}
整个量正比于 $q^2$，并且在 $q^2\to 0$ 极限下显然为零。这与轴矢量反常的预言形成对比。取 \eqref{eq:19.108} 右端的矩阵元，我们发现
\begin{equation}
  iq_\mu M^{\mu\nu\lambda} = \frac{e^2}{4\pi^2}\ \epsilon^{\nu\lambda\alpha\beta}
    p_\alpha k_\beta .
  \label{eq:19.113}
\end{equation}
如果出现在 \eqref{eq:19.112} 中的某个形状因子包含 $q^2$ 中的极点，这个冲突就可以被解决。这样的极点可以通过图~\ref{fig:19.8} 所示的过程产生，其中流产生一个随后衰变为两个光子的 $\pi^0$ 介子。流产生介子的振幅由 \eqref{eq:19.88} 给出。让我们把 $\pi$ 介子衰变振幅参数化为
\begin{equation}
  i\mathcal{M}(\pi^0 \to 2\gamma) = iA\,e\ \epsilon^{\nu\lambda\alpha\beta}\,
    p_\alpha k_\beta\ \epsilon_\nu(p)\epsilon_\lambda(k) ,
  \label{eq:19.114}
\end{equation}
其中 $A$ 是待定常数。则图~\ref{fig:19.8} 的过程对 \eqref{eq:19.109} 中定义的振幅 $M^{\mu\nu\lambda}$ 的贡献为
\begin{equation}
  \bigl( iq^\mu f_\pi \bigr)\ \frac{1}{q^2}\ \bigl( iA\,e\,
    \epsilon^{\nu\lambda\alpha\beta}p_\alpha k_\beta \bigr) .
  \label{eq:19.115}
\end{equation}
这是对形状因子 $M_1$ 的贡献，
\begin{equation}
  M_1 = \frac{f_\pi}{q^2}\cdot A\,e ,
  \label{eq:19.116}
\end{equation}
再加上在 $q^2 = 0$ 处正则的项。现在，通过令 \eqref{eq:19.112} 等于 \eqref{eq:19.113}，我们用反常的系数来确定 $A$：
\begin{equation}
  A = \frac{1}{4\pi^2 f_\pi} .
  \label{eq:19.117}
\end{equation}
由衰变矩阵元 \eqref{eq:19.114}，可以直接算出 $\pi^0$ 的衰变率。注意，尽管我们是在无质量 $\pi^0$ 的极限下算出衰变矩阵元的，我们必须提供依赖于 $\pi^0$ 质量的物理上正确的运动学。包括全同粒子相空间的因子 1/2，我们发现
\begin{equation}
  \Gamma(\pi^0 \to 2\gamma) = \frac{1}{2}\ \frac{1}{32\pi m_\pi}\ \int
    d\Omega\ |\mathcal{M}|^2 = \frac{\alpha^2 m_\pi^3}{64\pi^3 f_\pi^2} .
  \label{eq:19.118}
\end{equation}
这个公式与实验极好地一致。轴反常给出的 $\pi^0$ 寿命预言，与 $f_\pi$ 的测量值相结合，是对 Adler-Bell-Jackiw 反常的惊人确认。

\begin{figure}[t]
  \centering
  \includegraphics[width=0.5\textwidth]{fig19.8.png}
  \caption{导致轴矢量流形状因子 $M_1$ 中极点的贡献。}
  \label{fig:19.8}
\end{figure}

\section{手征反常与手征规范理论}

我们迄今讨论的反常出现在具有类矢量耦合的理论中，例如 QED 和 QCD。在这些理论中，反常只影响轴矢流，而轴矢流与规范对称性无关，因此理论的自洽性不会受到危害。然而，在\emph{手征规范理论}中，规范场以不同的方式耦合到左手与右手费米子。在这样的理论中，轴反常可以影响规范流本身的守恒，这将破坏理论的自洽性。我们现在将研究这个问题。

考虑一个规范群为 $G$ 的手征规范理论，其中左手费米子在 $G$ 的表示 $R$ 中变换。一个典型的例子是左手夸克和轻子的理论，弱相互作用规范群 $SU(2)$ 只耦合到左手双重态：
\begin{equation}
  \psi_L = \begin{pmatrix}\nu_L\\ e_L\end{pmatrix},\qquad
  q_L = \begin{pmatrix}u_L\\ d_L\end{pmatrix},
  \label{eq:19.124}
\end{equation}
而没有相应的右手耦合。

这样的理论的规范流接收来自图~\ref{fig:19.9} 所示三角图的修正。这些图是我们在第~19.2 节研究的反常图的类似物，但现在外流是规范流。规范流 $j^{\mu a}$ 散度中的反常由下式给出
\begin{equation}
  \partial_\mu j^{\mu a} = \frac{g^2}{32\pi^2}\ \epsilon^{\mu\nu\rho\sigma}\
    F^b_{\mu\nu}F^c_{\rho\sigma}\ A^{abc} ,
  \label{eq:19.131}
\end{equation}
其中
\begin{equation}
  A^{abc} = \mathrm{tr}\bigl[ t^a\{t^b, t^c\} \bigr] ,
  \label{eq:19.132}
\end{equation}
其中迹在费米子表示 $R$ 上求。

\begin{figure}[t]
  \centering
  \includegraphics[width=0.5\textwidth]{fig19.9.png}
  \caption{在手征规范理论中对规范对称性流的反常有贡献的图。}
  \label{fig:19.9}
\end{figure}

这个方程蕴含，除非 $A^{abc}$ 为零，否则流 $j^{\mu a}$ 不守恒。因子 \eqref{eq:19.132} 在 $(a, b, c)$ 中完全对称，所以这个条件与哪个流被当作外算符无关。正如我们在第~19.1 节和第~19.2 节中描述的，我们可以改变图的正则化以使外流守恒，但代价是违反图中另外两个流之一的守恒。

由于具有局域规范不变性的理论的整个构造都基于精确全局对称性的存在，$j^{\mu a}$ 守恒的违反会损害理论的结构。例如，图~\ref{fig:19.9} 形式的三角图现在将产生发散的规范玻色子质量项，并将破坏第~16 章讨论的三点和四点顶角之间精妙的关系。这些来自 Ward 恒等式的关系对于保证非物理态的抵消和 $S$ 矩阵的幺正性是必要的。避免这个问题的唯一方法是坚持 $A^{abc} = 0$ 作为手征规范理论的基本自洽性条件。%
\footnote{D.~J.~Gross and R.~Jackiw, \emph{Phys.\ Rev.\ D} \textbf{6}, 477
(1972)。} 满足这个条件的规范理论被称为\emph{无反常的}。

作为应用这个条件的例子，考虑我们在 \eqref{eq:19.124} 中给出的原型弱相互作用规范理论。如果图~\ref{fig:19.9} 中的两个规范玻色子是 $SU(2)$ 规范玻色子，且流 $j^{\mu a}$ 是 $SU(2)$ 规范流，我们将通过代入 $t^a = \tau^a = \sigma^a/2$ 并利用关系 $\{\sigma^b, \sigma^c\} = 2\delta^{bc}$ 来求值 \eqref{eq:19.132}。这给出
\begin{equation}
  A^{abc} = \mathrm{tr}\bigl[ \sigma^a\cdot 2\delta^{bc} \bigr] = 0 ,
  \label{eq:19.133}
\end{equation}
所以自洽性条件得到满足。如果 \eqref{eq:19.124} 中的费米子也耦合到电磁学，还存在一个额外的自洽性条件，我们将通过把图~\ref{fig:19.9} 中的流取为电磁流来找到它。这种情形下的因子 $A^{abc}$ 为
\begin{equation}
  \mathrm{tr}\bigl[ Q\{t^b, t^c\} \bigr] ,
  \label{eq:19.134}
\end{equation}
其中 $Q$ 是电荷矩阵。如果我们像在 \eqref{eq:19.133} 中那样化简，迹 \eqref{eq:19.134} 变为
\begin{equation}
  \frac{1}{2}\mathrm{tr}[Q]\,\delta^{bc} .
  \label{eq:19.135}
\end{equation}
这个因子正比于费米子电荷之和，它无论对夸克还是对轻子都不为零。然而，如果我们对一个夸克双重态和一个轻子双重态求和，色因子为 3，我们发现
\begin{equation}
  \mathrm{tr}[Q] = 3\cdot\Bigl( \frac{2}{3} - \frac{1}{3} \Bigr)
    + (0 - 1) = 0 .
  \label{eq:19.136}
\end{equation}
令人惊奇的是，只有当理论包含数目相等的夸克和轻子双重态时，\eqref{eq:19.124} 所描述的弱相互作用规范理论才能与 QED 自洽地组合。

我们通过更一般地导出手征规范理论无反常的条件来完成本节。我们将首先推导反常因子 $A^{abc}$ 的一些基本性质，然后把它们应用到具有简单规范群的手征规范理论。

如果费米子表示 $R$ 是实的，$R$ 等价于其共轭表示 $\bar{R}$。因此，正如我们在 \eqref{eq:15.82} 下方描述的，$t^a_R$ 通过幺正变换与 $t^a_{\bar{R}} = -(t^a_R)^T$ 相关联。由于 \eqref{eq:19.132} 对 $t^a$ 的幺正变换不变，我们可以把 $t^a_R$ 替换为 $t^a_{\bar{R}}$。则
\begin{equation}
  \begin{aligned}
    A^{abc} &= \mathrm{tr}\Bigl[ (-t^a)^T\ \bigl\{ (-t^b)^T, (-t^c)^T
      \bigr\} \Bigr]\\
    &= -\mathrm{tr}\Bigl[ \{ t^c, t^b \}\ t^a \Bigr] = -A^{abc} .
  \end{aligned}
  \label{eq:19.137}
\end{equation}
因此，如果 $R$ 是实的，规范理论自动无反常。作为一个特例，任何 Dirac 费米子的规范理论都是无反常的。

在更一般的情形下，我们可以通过注意到 $A^{abc}$ 是规范群 $G$ 的一个在伴随表示中具有三个指标且完全对称的不变量来简化它的计算。对某些可能的群，合适的不变量可能不存在，在那些情形下 $A^{abc}$ 必为零。例如，在 $SU(2)$ 中，伴随表示具有自旋 1。两个自旋 1 多重态的对称乘积给出自旋 0 加自旋 2，没有自旋 1 分量。因此，不存在耦合两个自旋 1 指标以给出自旋 1 的对称张量。于是因子 \eqref{eq:19.132} 在任何 $SU(2)$ 规范理论中都必为零。我们在式~\eqref{eq:19.133} 的一个显式例子中看到这一点发生。

在 $SU(n)$ 群中，$n > 3$，存在所要求类型的唯一对称不变量 $d^{abc}$。它出现在基本表示的表示矩阵的反对易子中：
\begin{equation}
  \{ t^a_n, t^b_n \} = \frac{1}{n}\delta^{ab} + d^{abc}t^c_n .
  \label{eq:19.138}
\end{equation}
这个不变量的唯一性蕴含，在 $SU(n)$ 规范理论中，任何形如 \eqref{eq:19.132} 的迹都正比于 $d^{abc}$。对每个表示 $r$，我们可以定义一个反常系数 $A(r)$：
\begin{equation}
  \mathrm{tr}\bigl[ t^a_r\{t^b_r, t^c_r\} \bigr] = \frac{1}{2}A(r)\,d^{abc}
  .
  \label{eq:19.139}
\end{equation}
对于基本表示，我们可以从 \eqref{eq:19.138} 看到 $A(n) = 1$。由 \eqref{eq:19.137} 的论证推出
\begin{equation}
  A(\bar{r}) = -A(r) .
  \label{eq:19.140}
\end{equation}
对于更高阶的表示，反常系数可以用类似于我们在第~15.4 节中用来计算 $C_2(r)$ 的方法导出。例如，我们在习题~19.3 中证明，如果 $a$ 和 $s$ 是对应于反对称和对称二指标张量的 $SU(n)$ 表示，则
\begin{equation}
  A(a) = n - 4,\qquad A(s) = n + 4 .
  \label{eq:19.141}
\end{equation}
如果费米子多重态 $R$ 的各个不可约分量的反常系数之和为零，则 $SU(n)$ 规范理论是无反常的。例如，具有表示内容
\begin{equation}
  R = a + (n-4)\bar{n}
  \label{eq:19.142}
\end{equation}
的左手费米子 $SU(n)$ 规范理论是无反常的。

在 \eqref{eq:15.72} 下方列出的各种简单 Lie 群中，只有 $SU(n)$、$SO(4n+2)$ 和 $E_6$ 具有复表示。其中，只有 $SU(n)$ 和 $SO(6)$（具有与 $SU(4)$ 相同的 Lie 代数）具有构造反常所要求类型的对称不变量。基于 $SO(4n+2)$（$n > 2$）和 $E_6$ 的规范理论自动无反常。群 $SO(10)$ 和 $E_6$ 已被提议为粒子物理大统一规范对称性的候选者，我们将在第~22.2 节讨论它。

对手征规范理论的表示内容还有一个进一步的约束，它来自考虑其与引力的耦合。可以证明，图~\ref{fig:19.9} 的图在用规范流 $j^{\mu a}$ 和外引力场计算时给出反常贡献。乘在这个图上的群论因子为
\begin{equation}
  \mathrm{tr}\bigl[ t^a \bigr] .
  \label{eq:19.143}
\end{equation}
如果规范群是非阿贝尔的，这个因子自动为零。然而，如果理论的规范群包含 $U(1)$ 因子，除非每个 $U(1)$ 生成元都是无迹的，否则理论不能与引力自洽地耦合。%
\footnote{L.~Alvarez-Gaum\'e and E.~Witten, \emph{Nucl.\ Phys.\ B} \textbf{234},
269 (1984)。}

一旦我们构造了自洽的手征规范理论，我们还面临一个额外的问题：找到在这个理论中自洽计算的方案。在类矢量规范理论中，我们可以用维数正则化来定义紫外发散图。这保证发散图将以尊重局域规范不变性 Ward 恒等式的方式被正则化。要把维数正则化推广到手征规范理论，我们需要引入 $\gamma^5$ 的维数延拓。在第~19.2 节中用于定义手征流的 't Hooft-Veltman 定义并不令人满意，因为这个定义不明显尊重规范流的守恒。一个有用的替代程序是把 $\gamma^5$ 形式上定义为与所有 $\gamma^\mu$ 反对易的对象。这个方案对不正比于 $\epsilon^{\mu\nu\rho\sigma}$ 的振幅给出不含糊的、规范不变的结果，至少到两圈阶。在第~21.3 节中，我们将使用这个方案来计算弱相互作用理论中的圈图。作为最后的手段，人们总可以用非规范不变的调节器来计算，并向理论添加非规范不变的抵消项，使得规范理论 Ward 恒等式保持有效。

\section{标度不变性的反常破缺}

还有一个重要的对称性例子：它在经典层面是不变性，却被量子修正破缺。这就是具有无量纲耦合常数的无质量场论的经典\emph{标度不变性}。在第~12 章中，我们看到没有经典有量纲参数的量子场论仍然通过紫外发散的（正则化）或者等价地通过耦合常数的跑动依赖于一个质量标度。我们已经看到如何用 Callan-Symanzik 方程分析这种对重整化标度的诱导依赖。在本节中，我们将证明量子修正对经典标度不变性的违反如何可以被描述为流守恒反常。

在本书中，我们一直避免对量子场论的能量-动量张量做仔细处理。在第~2.2 节中，我们用 Noether 定理证明了量子场论在时空平移下的不变性蕴含守恒张量 $T^{\mu\nu}$ 的存在。在第~9.6 节中，我们用泛函积分形式给出了这个结果的另一种推导。然而，为了讨论标度不变性理论，我们将需要能量-动量张量的一些更详细的性质。我们现在将直接陈述这些性质，并请读者参阅别处它们的推导。%
\footnote{接下来三段中给出的结论在 C.~G.~Callan, S.~Coleman, and R.~Jackiw,
\emph{Ann.\ Phys.} \textbf{59}, 42 (1970) 中被仔细地推导。}

由表达式 \eqref{eq:2.17} 和 \eqref{eq:9.99} 定义的张量 $T^{\mu\nu}$ 被称为\emph{正则能量-动量张量}。定义这个张量的表达式并不蕴含 $T^{\mu\nu}$ 是对称的。事实上，这个张量不必是对称的，而且在规范理论中，它不必是规范不变的。然而，总可以通过添加一个项把 $T^{\mu\nu}$ 转化为对称且规范不变的张量 $\Theta^{\mu\nu}$：
\begin{equation}
  \Theta^{\mu\nu} = T^{\mu\nu} + \partial_\lambda K^{\lambda\mu\nu} ,
  \label{eq:19.144}
\end{equation}
其中 $K^{\lambda\mu\nu}$ 在 $\mu$ 与 $\lambda$ 互换下反对称。所添加项的形式蕴含：如果 $T^{\mu\nu}$ 守恒，则 $\Theta^{\mu\nu}$ 守恒，并且全局能量-动量四矢量不变。

标量场论的标度变换可以定义为变量变换
\begin{equation}
  \phi(x) \to e^{-D\alpha}\phi(x e^{-\alpha}) ,
  \label{eq:19.146}
\end{equation}
其中 $D = 1$ 是场的正则质量量纲。标度变换在费米子和规范场理论中以类似方式定义。如果这个变换是经典拉格朗日量的不变性（在没有有量纲耦合时它将是如此），这个理论将拥有一个守恒流 $D^\mu$，称为\emph{伸缩流}。伸缩流与对称能量-动量张量 $\Theta^{\mu\nu}$ 有简单的关系：$D^\mu = x_\nu\Theta^{\mu\nu}$，从而
\begin{equation}
  \partial_\mu D^\mu = \Theta^\mu_{\ \mu} .
  \label{eq:19.147}
\end{equation}
从 Noether 定理推导这些结果并不直接。有一个更简单的推导，然而它使用的形式体系超出了本书的范围。如果所考虑的量子场论耦合到引力，则能量-动量张量可以被认定为引力场的源。这个能量-动量张量可以通过对物质场的拉格朗日量 $\mathcal{L}$ 关于时空度规 $g_{\mu\nu}(x)$ 变分来找到。这种构造给出一个显然对称且规范不变的张量，它结果正是
\begin{equation}
  \Theta^{\mu\nu} = \frac{2}{\sqrt{-g}}\ \frac{\delta\mathcal{L}}
    {\delta g_{\mu\nu}} .
  \label{eq:19.148}
\end{equation}
标度变换可以表示为时空度规的改变
\begin{equation}
  g_{\mu\nu}(x) \to e^{2\alpha}g_{\mu\nu}(x) .
  \label{eq:19.149}
\end{equation}
把 \eqref{eq:19.148} 和 \eqref{eq:19.149} 结合起来，我们看到这个变换引起的拉格朗日量改变正好是 $\Theta^{\mu\nu}$ 的迹。根据 Noether 定理，这将等于相应流的散度，使我们回到式~\eqref{eq:19.147}。

在 QED 中，不难猜出对称能量-动量张量的形式：
\begin{equation}
  \Theta^{\mu\nu} = -F^{\mu\lambda}F^{\nu}_{\ \lambda}
    + \frac{1}{4}g^{\mu\nu}(F_{\lambda\sigma})^2
    + \bar{\psi}\gamma^{(\mu}iD^{\nu)}\psi
    - g^{\mu\nu}\bar{\psi}(i\slashed{D} - m)\psi .
  \label{eq:19.150}
\end{equation}
这是一个规范不变的对称张量，它导致总能量的熟悉表达式，
\begin{equation}
  H = \int d^3x\ \Bigl\{ \frac{1}{2}(\mathbf{E}^2 + \mathbf{B}^2)
    + \psi^{\dagger}\bigl( -i\bm{\alpha}\cdot\nabla + m\beta \bigr)\psi
    \Bigr\} .
  \label{eq:19.151}
\end{equation}
为了将来参考，我们注意到这些结果在任何时空维数 $d$ 下在经典层面都成立。在四维中，规范场项的迹自动抵消。在利用作为算符运动方程成立的 Dirac 方程之后，我们发现 $\Theta^{\mu\nu}$ 的迹由下式给出
\begin{equation}
  \Theta^\mu_{\ \mu} = m\bar{\psi}\psi ,
  \label{eq:19.152}
\end{equation}
并且经典地，如果 $m = 0$，它确实为零。

当包含量子修正时，我们知道标度变换不是理论的对称性，因为参照更大标度的同一个理论含有重整化耦合常数的不同值。重整化耦合的移动为
\begin{equation}
  g \to g + \alpha\beta(g) ,
  \label{eq:19.153}
\end{equation}
并且拉格朗日量的相应改变为
\begin{equation}
  \delta\mathcal{L} = \alpha\beta(g)\,\frac{\partial\mathcal{L}}
    {\partial g} .
  \label{eq:19.154}
\end{equation}
因此，当包含量子修正时，经典标度不变理论中伸缩流的方程应当读作
\begin{equation}
  \partial_\mu D^\mu = \Theta^\mu_{\ \mu} = \beta(g)\,\frac{\partial
    \mathcal{L}}{\partial g} .
  \label{eq:19.155}
\end{equation}
在无质量 QED 中，我们可以通过重新标度规范场使耦合常数 $e$ 从协变导数中移除，从而最有用地写出这个公式：$eA_\mu\to A_\mu$。则 $e$ 只出现在项中
\begin{equation}
  \mathcal{L} \supset -\frac{1}{4e^2}(F_{\lambda\sigma})^2 + \cdots ,
  \label{eq:19.156}
\end{equation}
并且式~\eqref{eq:19.155} 读作
\begin{equation}
  \Theta^\mu_{\ \mu} = \frac{\beta(e)}{2e}\ \frac{1}{2}(F_{\lambda\sigma})^2
  = \frac{e^2}{24\pi^2}\ \frac{1}{4}(F_{\lambda\sigma})^2 .
  \label{eq:19.157}
\end{equation}
这个关系表明对称能量-动量张量的迹作为量子修正的结果取非零值，它被称为\emph{迹反常}。

我们应该能够在一个计算中直接检验迹反常方程 \eqref{eq:19.157}。迹反常的物理内容是：经典地为零的能量-动量张量的迹，从耦合常数的跑动接收到非零贡献。这个反常具有重要的物理推论，例如在强子的质量谱和 QCD 的能量-动量张量结构中。

\problems

\textbf{19.1} \emph{四维中的费米子数不守恒。} 在这个问题中，我们研究电场与磁场平行的背景场位形中费米子数的不守恒。考虑背景场中 $\mathbf{E}$ 与 $\mathbf{B}$ 平行的四维无质量 QED。证明反常方程导致费米子数的违反，并用场强计算费米子数违反的速率。

\textbf{19.2} 在正文中，我们用两种方法推导了四维 QED 中的轴矢量反常：算符乘积方法和泛函积分方法。在这个问题中，直接从三角图计算反常，并验证结果是 \eqref{eq:19.45}。你将需要仔细地正则化三角图，使用维数正则化，并跟踪 $\gamma^5$ 约定。

\textbf{19.3} \emph{更高表示的反常系数。} 在正文中，我们陈述了 $SU(n)$ 的反对称和对称二指标张量表示的反常系数 $A(a) = n-4$ 和 $A(s) = n+4$。用类似于第~15.4 节中用来计算 Casimir 算符的方法推导这些结果。

\textbf{19.4} \emph{QCD 中的迹反常。} 按照第~19.5 节的分析，推导无质量 QCD 的迹反常。证明 QCD 的能量-动量张量的迹正比于 $\beta(g)\,G^a_{\mu\nu}G^{a\mu\nu}$，并讨论物理推论。

\textbf{19.5} 在正文中我们证明，耦合到引力的手征规范理论的反常因子为 $\mathrm{tr}[t^a]$。证明除非 $U(1)$ 电荷是无迹的，否则 $U(1)$ 规范理论不能与引力自洽地耦合，并讨论这对理论的物质内容施加的约束。
